\documentclass{ou}
\begin{document}
\thispagestyle{empty}
\begin{center}
\shortstack{\FGAVO}\vspace{4\baselineskip}
\textbf{РЕФЕРАТ}\\
\setstretch{2}
Информационные системы и~технологии\\
\setstretch{1}
Мультимедиа\vspace{2\baselineskip}\end{center}
\begin{flushleft}
\doublespacing
\begin{tabularx}{\textwidth}{@{} p{0.4\textwidth} l l @{}}
Руководитель & ст. преподаватель & И. Р. Нигматуллин\\
Студенты     &                   &                  \\
\multirow{-2.27}{\hsize}{\singlespacing Современное понимание мультимедиа как~формы интеграции данных разных типов}
             & ГФ25-01Б, 152537684 & О. И. Бочанова\\
             & ГФ25-01Б, 152539986 & М. К. Гинтер\\
             & ГФ25-01Б, 152541471 & А. О. Безруких\\
             & ГФ25-01Б, 152536796 & А. В. Корюкова\\
\multirow{-2.27}{\hsize}{\singlespacing Основные виды мультимедийных объектов и~особенности их взаимодействия}
             & ГФ25-01Б, 152537671 & А. Е. Корикова\\
             & ГФ25-01Б, 152536528 & Д. А. Васильева\\
             & ГФ25-01Б, 152540016 & И. Н. Морозова\\
\multirow{-2.27}{\hsize}{\singlespacing Инструменты для~создания и~редактирования мультимедийного контента}
             & ГФ25-01Б, 152541294 & Э. Р. Махмудов\\
             & ГФ25-01Б, 152539829 & В. Д. Бобкова\\
\end{tabularx}
\end{flushleft}
\vspace*{\fill}
\begin{center}
\city
\end{center}
\newpage
\thispagestyle{empty}
\setstretch{1}
Продолжение титульного~листа реферата по~теме «Мультимедиа»
\begin{flushleft}
\doublespacing
\begin{tabularx}{\textwidth}{@{} p{0.4\textwidth} l l @{}}
Студенты     &                     &         \\
\multirow{-2.27}{\hsize}{\singlespacing Роль кодеков и~форматов в~обеспечении совместимости мультимедиа}
             & ГФ25-01Б, 152540765 & А. В. Гайворонская\\
             & ГФ25-01Б, 152540753 & М. И. Бобылев\\
             & ГФ25-01Б, 152539850 & Е. А. Аксаментова\\
\multirow{-2.27}{\hsize}{\singlespacing Применение мультимедиа в~образовательных и~коммуникационных процессах}
             & ГФ25-01Б, 152536115 & С. В. Идрисова\\
             & ГФ25-01Б, 152540131 & Е. Кичко\\
             & ГФ25-01Б, 152539526 & К. Н. Аркадьева\\
             & ГФ25-01Б, 152536056 & В. А. Карасова\\
\multirow{-2.27}{\hsize}{\singlespacing Правовые и этические аспекты работы с~мультимедийными материалами}
             & ГФ25-01Б, 152540046 & А. Д. Балцевич\\
             & ГФ25-01Б, 152537659 & В. Е. Зайцева\\
             & ГФ25-01Б, 152536592 & А. Ю. Бирюкова\\
            
\end{tabularx}
\end{flushleft}
\newpage
{
\sloppy
\tableofcontents
}
\newpage
\onehalfspacing
\section{Современное понимание мультимедиа как формы интеграции данных разных типов}
Под мультимедиа понимают не просто совокупность различных типов контента -- текст, звук, изображения, видео -- а комплексная концепция и технология, позволяющая объединять, обрабатывать и воспроизводить информацию разных типов в единой среде. Это объединение делает восприятие информации более гибким и многогранным.
Современное понимание мультимедиа -- это способ организации, структурирования, передачи и использования информации
\subsection{Определение и сущность мультимедиа}
Термин Мультимедиа изначально означает «многосредность» -- то есть использование нескольких медиа одновременно.

В более широком смысле мультимедиа -- это «объект-контейнер», в котором могут быть объединены разнообразные формы представления информации, а также -- средства для их обработки и воспроизведения. Так, интерактивные электронные учебники включают текстовые статьи, иллюстрации, аудиообъяснения и встроенные задания.

Мультимедиа -- согласованная система: аппаратные, программные компоненты, которые обеспечивают возможность хранения, обработки, передачи и воспроизведения мультимедийных данных. 

Например, стандарт MPEG-4 иллюстрирует, как мультимедиа может представляться в виде объектно-ориентированного формата, когда медийные объекты (видео, аудио, изображения, 3D-объекты, текст) интегрируются в~единую структуру. 

Таким образом, мультимедиа представляет собой новую форму информации -- не просто текстовая, звуковая или визуальная, а интегрированная, смешанная, многомерная.
\subsection{Аппаратно-программные средства}
Чтобы мультимедиа функционировала как единая система, необходимы специальные средства: аудио- и видеокарты, мультимедийные компьютеры, воспроизведение, программное обеспечение, обеспечивающее обработку, хранение, и взаимодействие с медийными объектами.  

Такие средства позволяют реализовать все этапы работы с информацией: ввод, обработку, хранение, передачу и вывод -- вне зависимости от типа данных.
\subsection{Интерактивность и линейность}
Одной из ключевых особенностей современного мультимедиа является возможность интерактивного взаимодействия. В отличие от~традиционных, «линейных» форм (например, видеофильм или аудиозапись, где пользователь потребляет контент в заранее заданной последовательности, без вмешательства), мультимедиа может предоставлять пользователю контроль -- нелинейное мультимедиа (интерактивное, гипермедиа) даёт свободу: пользователь может выбирать, что и когда просматривать, переходить между разделами, управлять порядком потребления информации.

Интерактивность ценна в образовании, тренингах, презентациях: она позволяет подстраивать материал под нужды пользователя, вовлечение, самостоятельный выбор пути познания. 

Благодаря интерактивности мультимедиа перестаёт быть пассивным носителем -- оно становится средством коммуникации, диалога, исследования. Пользователь не просто потребляет, но участвует, влияет, выбирает.
\subsection{Образование и обучение}
Одной из ключевых сфер применения мультимедиа остаётся образование. Мультимедиа позволяет создавать электронные учебные ресурсы, тренажёры, интерактивные курсы, видео- и аудиоуроки, что делает процесс обучения более гибким, индивидуализированным, продуктивным. 

Применение мультимедиа на уроках помогает удержать внимание, сделать материал более понятным и наглядным, что особенно важно при~сложных, абстрактных концепциях. 
\subsection{Культура, медиа, коммуникации}
В более широком социально-культурном контексте мультимедиа трансформирует способ, которым мы создаём, распространяем и воспринимаем информацию. 

Через мультимедиа художники, дизайнеры, создатели медийного контента могут реализовывать новые формы творчества, сочетая разные средства -- звук, изображение, текст, интерактивность -- и создавая гибридные формы искусства. 

В медийной сфере мультимедиа способствует появлению новых форм контента -- мультимедийных публикаций, трансмедийных проектов, мультимедийной журналистики, интерактивных презентаций, digital-артов. Это расширяет возможности коммуникации, делает её более динамичной, адаптивной и ориентированной на восприятие человека в цифровую эпоху. 

Мультимедиа, как современное понятие и технология, -- это гораздо больше, чем просто сочетание текста, звука, изображения и видео. Это интеграция разных типов данных, их представления, синтеза, структурирования таким образом, чтобы информация стала более доступной, наглядной, гибкой, адаптивной и глубокой. А интерактивность и нелинейность, заложенные в мультимедиа, превращают пользователя из пассивного потребителя в~активного участника, давая возможность выбирать и взаимодействовать.
\newpage
\section{Основные виды мультимедийных объектов и~особенности их взаимодействия}
Мультимедиа сегодня становится неотъемлемой составляющей нашей жизни. Благодаря мультимедиа мир информационных технологий преобразился совершенно чудесным образом: он стал ярким, увлекательным и легко воспринимаемым. Системы мультимедиа после своего появления произвели поистине революционные изменения практически во всех областях деятельности человека: в образовании, компьютерном обучении и тренингах, во многих сферах профессиональной деятельности, науки, искусства, в сфере развлечений (компьютерные игры, видео, аудио) и т. д\cite{LPGenerator}.

Появление систем мультимедиа обусловлено как требованиям практики, так и развитием теории в области информационных технологий. Однако, резкий рывок, произошедший в этом направлении за последние несколько лет, обеспечен прежде всего развитием технических и системных средств. Это и прогресс в развитии персональной электронно-вычислительной машиной (ПЭВМ): резко возросшие объемы оперативной и внешней памяти, быстродействие, графические возможности, и достижения в области «некомпьютерной» видеотехники, лазерных дисков, а также их массовое внедрение. Важную роль сыграла также разработка методов быстрого и~эффективного сжатия данных \cite{Katsko_2015}.

Современные персональные компьютеры в массе своей предоставляют аппаратную основу для большинства классических применений мультимедиа. Современные компьютерные приложения профессионального уровня стремятся как можно быстрее и увереннее завоевать популярность среди пользователей, привлечь их внимание, занять соответствующую часть рынка. И~элементы мультимедиа очень этому способствуют. Поэтому изучение вопросов, связанных с мультимедийными технологиями, системами, средствами разработки чрезвычайно важны и актуальны для любых специалистов компьютерных отраслей\cite{digitalocean}.
\subsection{Текст}
Текст -- наиважнейший элемент сферы информации. В мультимедиа невозможно обойтись без текста. Без текста придется использовать большое количество изображений, чтобы передать требуемые мысли. Звуковое сопровождение и голос комментатора, конечно, помогут пользователям, но это в скором времени станет утомительным. Восприятие звуков и речи требует гораздо больше внимания и сил, чем чтение текста.

Текст становится важнейшим средством передачи информации, и он не утратил своего значения до сегодняшнего дня. Появление компьютеров подняло работу с текстовой информацией на недостижимую ранее высоту.

С появлением Интернета появляется новый тип текста -- гипертекст, использующийся в создании Web-страниц.

Гипертекст -- это мощная система связанных слов и фраз, позволяющая легко перемещаться по Web-страницам. Она связывает фразу или слово одной страницы с любой другой страницей, абзацем, фразой или словом\cite{Panasenko_2022}.
\subsection{Видео}
Видео -- электронная технология формирования, записи, обработки, передачи, хранения и воспроизведения подвижного изображения, основанная на принципах телевидения, а также аудиовизуальное произведение, записанное на физическом носителе (видеокассете, видеодиске и т. п.).

Для работы с видео помимо обычных видеокарт, предназначенных для вывода графической информации на монитор компьютера (их уместнее называть графическими адаптерами), для работы с видео используются специальные видеокарты \cite{Katsko_2015}.
\subsection{Графика}
Компьютерная графика является важным информационным элементом.

Изображения обычно поступают в компьютер следующими основными способами:
\begin{enumerate}[leftmargin=0em, itemindent=3.75em, nosep, label=\arabic*)]
\item Вводятся через сканер;
\item Выбираются из файлов, содержащих набор графических вставок и~поставляемых специализированными фирмами;
\item Создаются заново пользователями с помощью пакетов графических программ.
\end{enumerate}

После этого изображения можно подвергнуть последующей обработке различными способами.

По способам задания изображений графику делятся на~категории:
\begin{itemize}[label=-, leftmargin=0em, itemindent=3.35em, nosep]
\item Двумерная графика;
\item Трёхмерная графика;
\end{itemize}

Двумерная графика классифицируется и по типу представления графической информации, и следующими из него алгоритмами обработки изображений. Обычно компьютерную графику разделяют на векторную и растровую, хотя обособляют ещё и фрактальный тип представления изображений. Они отличаются принципами формирования изображения при~отображении на экране монитора или при~печати на бумаге.

Растровая графика -- способ сохранения изображения, при котором изображение является матрицей элементов -- пикселов (представляется в виде набора окрашенных точек). Размер растровой картинки может быть задан как~Х пикселов по ширине и Y пикселов по высоте. 

Векторные изображения -- способ сохранения изображения, при~котором изображение сохраняется в виде геометрического описания объектов, составляющих рисунок. Простейший геометрический объект -- это вектор, поэтому векторные изображения есть набор математического описания векторов. Эти изображения могут также включать в себя данные в формате растровой графики (представления изображения виде совокупности отрезков и дуг и т. д.). 

Фрактальная графика основана на математических вычислениях. Фрактальные изображения применяются в самых разных сферах, начиная от~создания обычных текстур и фоновых изображений и заканчивая фантастическими ландшафтами для компьютерных игр или книжных иллюстраций. Создаются фрактальные изображения путем математических расчетов. Базовым элементом фрактальной графики является математическая формула -- это означает, что никаких объектов в памяти компьютера не~хранится, и изображение строится исключительно на основе уравнений. 

Трёхмерная графика оперирует с объектами в трёхмерном пространстве. Трехмерная графика бывает полигональной и воксельной. Трёхмерная компьютерная графика широко используется в кино, компьютерных играх\cite{Panasenko_2022}.
\subsection{Звук}
Звук является наиболее выразительным элементом мультимедиа.

Звук -- это колебания воздуха или любой другой среды, в~которой~он распространяется. Звук характеризуется амплитудой (силой) и частотой (количеством колебаний в секунду).

Правильное использование звука определяет качество проекта, создает требуемое настроение. Это может быть как смысловая речь на любом языке, так и музыка, фоновое звучание и звуковые эффекты\cite{Katsko_2015}.
\subsection{Анимация}
Анимация -- искусственное представление движения путем отображения последовательности рисунков, кадров с частотой, при которой обеспечивается целостное зрительное восприятие образов. Для плавного воспроизведения анимации необходима частота кадров, не менее 10 кадров в секунду.

Термин «анимировать» дословно означает «оживить» изображение. Теория анимации базируется на положении о способности человеческого глаза сохранять на сетчатой оболочке след увиденного и соединять быстро меняющиеся изображения в единый зрительный ряд. Это создает иллюзию непрерывного движения.

Виды анимаций:
\begin{enumerate}[leftmargin=0em, itemindent=3.75em, nosep, label=\arabic*)]
\item Классическая (традиционная) анимация;
\item Стоп-кадровая (кукольная) анимация;
\item Спрайтовая анимация;
\item Морфинг;
\item Цветовая анимация;
\item 3D-анимация;
\item Захват движения\cite{Panasenko_2022}.
\end{enumerate}
Классическая  анимация представляет собой поочередную смену рисунков, каждый из которых нарисован отдельно. Это очень трудоемкий процесс, так~как аниматорам приходится отдельно создавать каждый кадр.

При стоп-кадровой анимации размещенные в пространстве объекты фиксируются кадром, после чего их положение изменяется и вновь фиксируется.

Спрайтовая анимация реализуется при помощи языка программирования. Анимация, основанная на движущихся объектах, называется спрайтовой анимацией, а объекты именуются спрайтами. Анимация осуществляется перемещением спрайтов и их заменой. Для управления спрайтами обычно имеется язык программирования. Первый плоские игры были основаны на~спрайтах.

Морфинг -- преобразование одного объекта в другой за счет генерации заданного количества промежуточных кадров. Программное обеспечение морфинга генерирует заданное число промежуточных кадров, которое обеспечивает плавный переход начального образа в конечный.

Цветовая анимация -- при ней изменяется лишь цвет, а не положение объекта.

3D-анимация создается при помощи специальных программ (например, 3D MAX). Картинки получаются путем визуализации трехмерной сцены, а~каждая сцена представляет собой набор объектов, источников света, текстур. Трехмерная графика в отличие от двухмерной дает более реалистичное представление образов.

Захват движения -- первое направление анимации, которое дает возможность передавать естественные, реалистичные движения в реальном времени. Датчики прикрепляются на живого актера в тех местах, которые будут приведены в соответствие с контрольными точками компьютерной модели для ввода и оцифровки движения. Координаты актера и его ориентация в пространстве передаются графической станции, и анимационные модели оживают\cite{Katsko_2015}.
\subsection{Классификация по структуре: линейные и нелинейные мультимедийные системы}
Мультимедиа может быть разделена на линейную (без обратной связи) и~нелинейную (интерактивную) среду.

Линейное мультимедиа -- это тип мультимедийного продукта, на~события которого потребитель повлиять не может. После запуска линейный продукт работает только по заданному сценарию. Примеры -- потоковое видео без~записи и возможности перемотки, статичные картинки, тексты без~гиперссылок. Фильм -- также пример линейного мультимедиа. Он идет последовательно, от начала до конца. Зритель не может поменять героев местами, повлиять на сюжет или изменить озвучку.

Линейные мультимедиа подходят, когда надо донести информацию для~большой группы людей. Допустим, на тренингах и семинарах, где десятки слушателей.

Нелинейный способ представления информации позволяет человеку, программам участвовать в выводе информации, взаимодействуя каким-либо образом со средством отображения мультимедийных данных. Участие в~данном процессе двух и более сторон называется «интерактивностью». Такой способ взаимодействия человека и компьютера наиболее полным образом представлен в категориях компьютерных игр. Такой способ представления мультимедийных данных иногда называется «гипермедиа».

Самым ярким примером нелинейных мультимедийных технологий являются компьютерные игры, а также разнообразная обучающая литература, в которой человеку предоставляется выбор различных действий. Принцип действия нелинейного мультимедиа заключается в том, что человек, использующий нелинейные мультимедийные технологии, может напрямую участвовать в выводе информации. Это осуществляет благодаря его взаимодействию с определенными средствами отображения различных мультимедийных объектов. Большинство современных мультимедийных продуктов можно отнести к нелинейным\cite{digitalocean}. 
\subsection{Взаимодействие мультимедийных объектов}
На синтаксическом, или формальном, уровне регулируются технические и пространственно-временные аспекты совмещения. К ним относится, прежде всего, точная синхронизация аудиоряда и видеоряда, анимации и звуковых эффектов, текстовых титров и речевой дорожки, нарушение которой ведёт к когнитивному диссонансу у пользователя. Сюда же относится пространственная композиция -- расположение элементов в интерфейсе согласно законам визуального восприятия, и формальная совместимость, обеспечиваемая использованием соответствующих кодеков и форматов данных.

Семантический уровень отвечает за согласованность смыслового содержания компонентов. На этом уровне происходит не дублирование, а взаимодополнение информации: визуальный ряд конкретизируется и~интерпретируется текстовыми пояснениями, в то время как аудиодорожка формирует эмоциональный контекст и атмосферу, облегчая декодирование вербальных и визуальных сообщений. Таким образом, мультимедийная система позволяет осуществлять редукцию информационной избыточности, где каждый модуль вносит уникальный вклад в общее смысловое поле.

Сложным является прагматический уровень, ориентированный на~взаимодействие системы с пользователем. Его центральным понятием выступает интерактивность -- способность системы реагировать на действия пользователя, изменяя состав, порядок, свойства медиакомпонентов. Это преобразует реципиента из пассивного потребителя информации в активного соучастника коммуникационного процесса. Развитием этого принципа является адаптивность, то есть способность системы гибко подстраивать набор и представление мультимедийных объектов под параметры контекста, такие~как характеристики устройства воспроизведения или пропускная способность канала связи\cite{Panasenko_2022}.
\newpage
\section{Инструменты для~создания и~редактирования мультимедийного контента}
В современном цифровом мире мультимедийный контент (текст, изображения, аудио, видео, анимация) стал основой коммуникации, маркетинга, образования и развлечений. Эффективное создание и редактирование такого контента требует использования специализированных инструментов. Цель данного доклада -- классифицировать и охарактеризовать ключевые категории программного обеспечения для работы с мультимедиа.
\subsection{Редактирование и создание изображений (графический дизайн)}
Эта категория охватывает инструменты для работы с растровой и~векторной графикой, создания иллюстраций и макетов.

Профессиональные пакеты для растровая графики:
\begin{itemize}[label=-, leftmargin=0em, itemindent=3.35em, nosep]
\item Adobe Photoshop: индустриальный стандарт. Используется для ретуши фотографий, сложного монтажа, создания цифровых картин, веб-дизайна и~работы с текстурой. Обладает неограниченными возможностями, но требует обучения;
\item Affinity Photo: мощный конкурент Photoshop. Предлагает аналогичный функционал за единоразовую плату, а не по подписке. Высоко ценится за~скорость и точность.
\end{itemize}

Программы для векторной графики:
\begin{enumerate}[leftmargin=0em, itemindent=3.75em, nosep, label=\arabic*)]
\item Adobe Illustrator: лидер в создании векторной графики (логотипы, иконки, иллюстрации). Векторные изображения можно масштабировать без~потери качества;
\item CorelDRAW: еще один мощный пакет для векторного дизайна, популярный в полиграфии и дизайне макетов;
\item Figma: онлайн-инструмент, ставший стандартом для UI/UX-дизайна (создание интерфейсов сайтов и приложений). Позволяет командам работать над проектом одновременно.
\end{enumerate}

Доступные и бесплатные альтернативы:
\begin{itemize}[label=-, leftmargin=0em, itemindent=3.35em, nosep]
\item GIMP (GNU Image Manipulation Program): Бесплатный и мощный редактор с открытым исходным кодом. Подходит для многих задач, которые решает Photoshop;
\item Canva: онлайн-платформа, ориентированная на непрофессионалов. Содержит огромную библиотеку шаблонов для социальных сетей, презентаций, плакатов. Идеален для быстрого создания качественного визуального контента без глубоких знаний дизайна.
\end{itemize}
\subsection{Производство и редактирование видео}
Инструменты для монтажа видео, наложения эффектов, цветокоррекции и работы со звуком.

Профессиональные рабочие станции (NLE -- Non-Linear Editors):
\begin{enumerate}[leftmargin=0em, itemindent=3.75em, nosep, label=\arabic*)]
\item Adobe Premiere Pro: Отраслевой стандарт для монтажа видео. Интегрируется с другими продуктами Adobe (After Effects, Audition). Используется для монтажа фильмов, ТВ-шоу и онлайн-видео.
\item DaVinci Resolve: Уникальный инструмент, сочетающий в себе профессиональный монтаж, лучшую в индустрии цветокоррекцию, а также функции для работы со звуком и визуальными эффектами. Существует бесплатная версия с огромным функционалом.
\item Final Cut Pro: Эксклюзив для macOS. Известен оптимизацией, скоростью работы, интерфейсом. Популярен среди видеоблогеро, независимых режиссеров.
\end{enumerate}

Программа для визуальных эффектов (VFX) и моушн-дизайна -- Adobe After Effects: Стандарт для создания анимированной графики, визуальных эффектов и композитинга. Используется для создания титров, анимации логотипов, спецэффектов.

Решения для начинающих и быстрого монтажа:
\begin{itemize}[label=-, leftmargin=0em, itemindent=3.35em, nosep]
\item Adobe Premiere Rush: упрощенная версия Premiere Pro, предназначенная для создания контента для социальных сетей на компьютерах и мобильных устройствах;
\item CapCut: бесплатное и очень популярное приложение для монтажа с~огромной библиотекой шаблонов, эффектов и музыки. Идеален для создания коротких роликов для TikTok, Reels, YouTube Shorts;
\item iMovie: бесплатное и простое решение для пользователей Apple, позволяющее быстро собрать качественный видеоотчет или семейное видео.
\end{itemize}
\subsection{Создание и обработка аудио}
Инструменты для записи, сведения, мастеринга и создания музыки.

Цифровые аудио рабочие станции (DAW):
\begin{enumerate}[leftmargin=0em, itemindent=3.75em, nosep, label=\arabic*)]
\item Adobe Audition: профессиональный инструмент для записи и сведения звука, популярный среди подкастеров и звукорежиссеров;
\item Ableton Live: лидер среди музыкантов-электронщиков, для живых выступлений и создания музыки;
\item FL Studio, Logic Pro, Cubase: мощные DAW для аранжировки и~сведения музыки любого жанра.
\end{enumerate}

Бесплатные и специализированные инструменты:
\begin{itemize}[label=-, leftmargin=0em, itemindent=3.35em, nosep]
\item Audacity: бесплатный, простой в использовании аудиоредактор с открытым исходным кодом. Отлично подходит для базовой записи и~редактирования (обрезка, шумоподавление, нормализация);
\item GarageBand: бесплатная DAW от Apple для пользователей macOS и iOS, идеальная точка входа для начинающих музыкантов.
\end{itemize}
\subsection{Создание анимации и 3D-графики}
Инструменты для оживления статичных изображений и создания трехмерных миров.

2D-анимация:
\begin{enumerate}[leftmargin=0em, itemindent=3.75em, nosep, label=\arabic*)]
\item Adobe Animate: стандарт для создания векторной 2D-анимации, часто используется для веб-баннеров и интерактивного контента.
\item Toon Boom Harmony: профессиональное решение для создания полноценных анимационных фильмов и сериалов.
\item Blender: мощный бесплатный пакет, который также включает в себя инструменты для 2D-анимации.
\end{enumerate}

3D-моделирование, анимация и визуализация:
\begin{itemize}[label=-, leftmargin=0em, itemindent=3.35em, nosep]
\item Blender: полнофункциональный пакет для 3D-творчества с открытым исходным кодом. Включает моделирование, скульптинг, анимацию, рендеринг и композитинг. Конкурирует с дорогими коммерческими аналогами;
\item Autodesk Maya \verb|&| 3ds Max: индустриальные стандарты для создания голливудских визуальных эффектов, анимации и архитектурной визуализации;
\item Cinema 4D: известен своим дружелюбным интерфейсом и мощными инструментами для моушн-дизайна, часто интегрируется с After Effects.
\end{itemize}

Выбор инструмента для создания и редактирования мультимедийного контента зависит от конкретных задач, уровня навыков пользователя и бюджета. Существует решение для любого случая: от бесплатных и простых редакторов для новичков до комплексных профессиональных студий, способных производить контент кинематографического качества. Ключевой тенденцией является демократизация технологий, позволяющая все большему числу людей реализовывать свои творческие идеи без необходимости многолетнего обучения и огромных финансовых вложений.
\newpage
\section{Роль кодеков и~форматов в~обеспечении совместимости мультимедиа}
\subsection{Понятие кодеков и форматов мультимедиа}
Кодек -- алгоритм, который преобразует мультимедийную информацию в удобное для хранения или передачи состояние и возвращает её обратно для~воспроизведения.

Каждый кодек выполняет две основные функции:
\begin{itemize}[label=-, leftmargin=0em, itemindent=3.35em, nosep]
\item Кодирование -- сжатие исходных данных;
\item Декодирование -- восстановление данных при воспроизведении.
\end{itemize}

Кодеки различаются по степени сжатия, качеству, скорости работы и~оборудованию, которое требуется для работы.

Формат или контейнер -- это структура файла, которая объединяет потоковые данные, созданные кодеками, и дополнительные элементы:
\begin{itemize}[label=-, leftmargin=0em, itemindent=3.35em, nosep]
\item Субтитры;
\item Метаданные;
\item Главы;
\item Несколько дорожек аудио;
\item Сведения о кодеке.
\end{itemize}
\subsection{Роль кодеков в обеспечении совместимости}
Если бы не существовало общепринятых кодеков, каждый производитель записывал бы видео и аудио «по-своему». Это означало бы, что:
\begin{enumerate}[leftmargin=0em, itemindent=3.75em, nosep, label=\arabic*)]
\item Видеоролик, снятый на один смартфон, не открывался бы на другом;
\item Телевизоры не умели бы воспроизводить большинство файлов;
\item Видеоредакторы требовали бы собственные форматы.
\end{enumerate}

Стандартизированные кодеки устраняют эти проблемы: они дают возможность любой программе или устройству понимать, как устроен медиапоток.

Большинство кодеков используют сжатие с потерями, поскольку это значительно уменьшает размер файлов. Например:
\begin{enumerate}[leftmargin=0em, itemindent=3.75em, nosep, label=\arabic*)]
\item Исходный поток 4K-видео может весить десятки гигабайт;
\item После кодирования H.265 тот же ролик весит в 5–10 раз меньше, сохраняя хорошее качество.
\end{enumerate}

Оптимизация важна и для потокового вещания. Чем лучше кодек сжимает данные, тем:
\begin{itemize}[label=-, leftmargin=0em, itemindent=3.35em, nosep]
\item Меньше требуется пропускная способность сети;
\item Стабильнее работает стриминг;
\item Быстрее загружается видео даже на слабых устройствах.
\end{itemize}
\subsection{Роль форматов (контейнеров) в совместимости мультимедиа}
Контейнер отвечает за то, как именно собраны воедино все элементы мультимедиа.

Для совместимости очень важно, чтобы устройства и программы могли:
\begin{itemize}[label=-, leftmargin=0em, itemindent=3.35em, nosep]
\item Распознать контейнер;
\item Определить, какие кодеки внутри;
\item Правильно прочитать структуру файла.
\end{itemize}

Например, AVI -- старый контейнер: поддерживает ограниченное число кодеков. MP4 -- современный и гибкий, поддерживает множество типов видео и аудио.

Фильм в контейнере MKV может содержать:
\begin{enumerate}[leftmargin=0em, itemindent=3.75em, nosep, label=\arabic*)]
\item Видео H.264;
\item Аудио DTS;
\item Русские и английские дорожки;
\item Набор субтитров;
\item Главы.
\end{enumerate}

Но если телевизор не поддерживает, например, DTS, он воспроизведёт видео, но без звука. Это классический пример несовместимости.
\subsection{Популярные кодеки и их применение}
Видео-кодеки. H.264:
\begin{itemize}[label=-, leftmargin=0em, itemindent=3.35em, nosep]
\item Самый распространённый в мире;
\item Используется YouTube, Zoom, Skype, видеокамеры;
\item Хорошо сжимает видео при высокой совместимости.
\end{itemize}

H.265:
\begin{enumerate}[leftmargin=0em, itemindent=3.75em, nosep, label=\arabic*)]
\item Улучшенная версия H.264 -- на 40–50\verb|%| эффективнее;
\item Требует более мощного оборудования;
\item Применяется в 4K/8K-контенте, Apple, Netflix.
\end{enumerate}

AV1:
\begin{itemize}[label=-, leftmargin=0em, itemindent=3.35em, nosep]
\item Свободный и бесплатный кодек нового поколения;
\item Создан Google, Netflix, Intel;
\item На 30\verb|%| эффективнее H.265;
\item Активно внедряется в YouTube и видеопотоки.
\end{itemize}

VP9:
\begin{enumerate}[leftmargin=0em, itemindent=3.75em, nosep, label=\arabic*)]
\item Кодек от Google для YouTube и WebM;
\item Используется в браузерах, поддерживающих HTML5-видео.
\end{enumerate}

Аудио-кодеки. MP3:
\begin{itemize}[label=-, leftmargin=0em, itemindent=3.35em, nosep]
\item Стандарт для музыки более 20 лет;
\item Сильно сжимает файлы, но снижает качество.
\end{itemize}

AAC:
\begin{enumerate}[leftmargin=0em, itemindent=3.75em, nosep, label=\arabic*)]
\item Популярнее MP3 в современных сервисах;
\item Используется Apple Music, YouTube.
\end{enumerate}

FLAC:
\begin{itemize}[label=-, leftmargin=0em, itemindent=3.35em, nosep]
\item Сжатие без потерь;
\item Используется меломанами и студиями.
\end{itemize}

Opus:
\begin{enumerate}[leftmargin=0em, itemindent=3.75em, nosep, label=\arabic*)]
\item Универсален для голоса и музыки;
\item Применяется в Discord, WhatsApp, WebRTC;
\item Отлично работает даже при слабом интернете.
\end{enumerate}
\subsection{Реальные примеры совместимости и несовместимости}
Zoom, Skype и Teams используют кодеки, которые работают стабильно даже при низкой скорости: H.264 и специфические кодеки для голоса. Благодаря этому конференции открываются на Windows, macOS, Android, iOS и даже в~браузере.

iPhone записывает видео в HEVC, но некоторые устройства Android не~поддерживают его аппаратно. Поэтому видео может:
\begin{itemize}[label=-, leftmargin=0em, itemindent=3.35em, nosep]
\item Открываться с ошибками;
\item Не воспроизводиться;
\item Требовать конвертации в H.264.
\end{itemize}

Netflix хранит один и тот же фильм в разных кодеках:
\begin{enumerate}[leftmargin=0em, itemindent=3.75em, nosep, label=\arabic*)]
\item H.264 -- для старых устройств;
\item HEVC -- для современных;
\item AV1 -- для премиальных клиентов или 4K-контента.
\end{enumerate}

Это гарантирует, что видео работает у всех.
\subsection{Тенденции развития кодеков и форматов}
AV1 становится стандартом, так как не требует лицензий. Это снижает затраты производителей и улучшает совместимость.

Новые графические чипы и процессоры имеют встроенную поддержку кодеков. Это делает воспроизведение более стабильным и энергоэффективным.

4K и 8K требуют ещё более эффективных кодеков, что ускоряет развитие H.266

Новые форматы, такие как CMAF, оптимизированы для потокового вещания и снижают задержку.
\newpage
\section{Применение мультимедиа в~образовательных и~коммуникационных процессах}
Тенденции развития современной системы высшего образования неразрывно связаны с широким внедрением в учебный процесс различных форм, методов и средств активного обучения. Одной из ведущих тенденций информатизации общества является развитие мультимедийных технологий, их проникновение в различные сферы социальной жизни: производство, бизнес, науку, образование, массовую потребительскую культуру. Обеспечивая богатство содержания и формы, сочетание различных видов текстовой, графической, речевой, музыкальной, видео-, фото- информации и разнообразие способов их извлечения, эти технологии формируют мультимедийное восприятие мира. В образовательном контексте они открывают принципиально новые возможности для организации учебного процесса, повышения его эффективности и развития творческого потенциала обучающихся. Однако их успешное внедрение требует решения ряда организационных, методических и~технических задач. 
\subsection{Понятие и сущность}
Мультимедиа как современный вид представления информации включает в себя: текст (в устной и письменной формах), статические изображения (таблицы, графики, иллюстрации), звук, видео, анимацию и др. Применение современных способов обработки аудиовизуальной информации (удобная навигация, гибкая порционная подача, интерактивность), одновременное яркое воздействие на различные рецепторы обучаемого позволяет интенсифицировать процесс обучения и повысить эффективность усвоения учебной информации.

Такая универсальность позволяет формировать чувственный опыт как основу обучения, представляя информацию в максимально приближенном к~реальности виде. Мультимедиа воздействует на обучаемого через различные каналы восприятия -- слуховой, зрительный и моторный, а также создает определенные эмоциональные ощущения.

По данным Computer Technology Research Corporation, издательской и~исследовательской компании, люди запоминают лишь 20\verb|%| того, что слышат, и 30\verb|%| того, что видят. Однако они могут помнить лишь 50\verb|%| того, что видят и~слышат, и до 80\verb|%| того, что они одновременно смотрят, чувствуют и~выполняют. Именно поэтому мультимедиа является бесценным инструментом обучения и преподавания\cite{journals_stmjournals_2024}.
\subsection{Преимущества и недостатки использования мультимедиа в~образовании}
Использование мультимедийных средств в образовании даёт следующие возможности: 
\begin{enumerate}[leftmargin=0em, itemindent=3.75em, nosep, label=\arabic*)]
\item Повышение эффективности учебного процесса; 
\item Развитие личных качеств обучающихся; 
\item Развитие коммуникативных и социальных способностей обучающихся; 
\item Расширение возможностей индивидуализации и дифференциации открытого и дистанционного обучения каждой личности с помощью компьютерных технологий, электронных информационных ресурсов; 
\item Отношение к обучаемому как активно получающего знаний субъекту, признание его достоинств; 
\item Учёт личного ответа и индивидуальных особенностей обучаемого; 
\item Самостоятельное ведение учебной деятельности; 
\item Воспитание умений адаптации к быстро меняющимся социальным условиями и успешного выполнения профессиональных задач с использованием современных образовательных технологий.
\end{enumerate}

Существуют некоторые проблемы в применении мультимедиа в процессе образования. В частности:
\begin{enumerate}[leftmargin=0em, itemindent=3.75em, nosep, label=\arabic*)]
\item Отвлечение внимания. Сложные методы представления информации могут отвлечь внимание обучающихся от темы изучения. Информации большого объёма, различные другие приложения, используемые в средствах мультимедиа, могут отвлечь внимание студентов; 
\item Сложности в создании учебных материалов. Создание элементов мультимедиааудио, видео графических и других элементов сложнее, чем создание традиционных материалов в виде текстов; 
\item Требует больше времени; 
\item Проблемы, возникающие в подготовке и использовании средств программного обеспечения и техники. Для эффективного использования мультимедийных средств обучения технические средства должны быть в~исправном состоянии. Презентация мультимедийных материалов требует большего качества и широких возможностей, чем редактирование текстов и~простого представления их другим студентам; 
\item Трудности, возникающие при чтении информации через экраны компьютеров. Чтение компьютерного текста сложнее, чем чтение текста, напечатанного на бумаге. Информация большого объёма, чтение газетных и~книжных вариантов, требующих полного ознакомления более удобно.
\end{enumerate}
\subsection{Практическая реализация и инструментарий мультимедиа}
За последние десятилетия в системе вузовского и школьного образования создано огромное количество различных электронных учебных пособий (учебники, справочные системы, тесты, энциклопедии, обучающие программы и т.п.).

Для создания электронного учебного пособия используются самые разнообразные программные и аппаратные средства, их применение обычно требует серьезной технологической подготовленности преподавателя, а также отнимает достаточно много времени на его разработку и отладку. При~этом их эффективность и уместность нередко вызывают сомнение, тем более что практическое применение подобных средств в учебных целях обычно возможно лишь при индивидуальной работе за компьютером (например, в~компьютерном классе). На деле же большинство вузовских аудиторий укомплектовано в лучшем случае только мультимедийным компьютером и~проектором (иногда в комплекте с интерактивной доской).

Из вышесказанного можно сделать вывод: для широкого использования в образовательном процессе вуза необходимо применять не слишком сложное программное обеспечение, позволяющее даже не имеющему специального образования преподавателю создавать яркие наглядные электронные пособия с использованием широкого спектра мультимедиа как вида учебной информации.

Таким распространенным программным продуктом, не~требующим сложных аппаратных ресурсов (компьютерный класс и т.п.) является программа PowerPoint (общеизвестное мощное средство для создания мультимедийных презентаций).

Технология создания и применения наглядных пособий с использованием данной программы широко известна и обычно не вызывает затруднений у преподавателей и студентов, однако нередко можно отметить низкую эффективность подобных разработок, связанную с различными факторами.

Основные ошибки использования презентационных технологий связаны с проектированием и конструированием мультимедийных учебных пособий. Подготовка наглядных электронных пособий для проведения лекции, практического или лабораторного занятия требует тщательной продуманности как в системе подачи, структуре учебного материала, так и в формах представления мультимедийной информации.

При проектировании будущего занятия с использованием мультимедиа разработчик должен определить его конкретные цели и задачи: изучение нового материала, закрепление, повторение и т.д. И только, исходя из~конкретно поставленной задачи, необходимо выбрать ту форму электронного учебного пособия, которая будет нацелена на максимальную эффективность усвоения учебной информации.

Мультимедийное наглядное пособие является вторичным продуктом по~отношению к исходной учебной информации. Оно являет собой некоторую информационную модель, отображающую и замещающую исходную информацию в простой и удобной для познания форме. 

Очевидно, что при выборе модели важно исходить из целей обучения, а~не~из идентичности графической информации и реальной действительности.

Мультимедиа используется для усиления наглядности занятия, но~при~этом не следует забывать, что в большинстве случаев электронное пособие служит именно визуальной (или звуковой) иллюстрацией вербальной учебной информации. Наглядное мультимедийное пособие обычно не является основным источником информации на занятии, поэтому нет необходимости максимально насыщать презентацию текстом.

Текстовая информация на экране должна быть минимизирована по~объему, представлять основу, стержень занятия, при~этом повышаются требования по отбору, структурированию, систематизации учебных знаний. Письменный текст в ходе презентации должен не накладываться на речь преподавателя, не подменять ее, а служить визуальной поддержкой вербальной информации.

Продолжая тему визуализации, важно отметить значение цветовых соотношений графического оформления презентаций (оптимальное количество цветов, особенности воздействия на познавательную деятельность и т.д.).

Не менее важно значение звука, который может быть как средством создания определенной эмоциональной атмосферы занятия, средством активизации внимания, так и дополнительным источником учебной информации.

Наконец, необходимо отметить особое значение современных возможностей мультимедиа: видео, анимация и интерактивность. Если графика, звук, речь учителя традиционно воздействуют в основном на~визуальный и аудиальный типы восприятия информации, то движение динамика, возможности непосредственного управления учебными процессами позволяют успешно и эффективно воспринимать учебную информацию кинестетическому типу восприятия.
\subsection{Применение мультимедиа в профессиональной и массовой коммуникации}
Помимо образовательной сферы, мультимедийные технологии кардинально преобразовали современные коммуникационные процессы, став их неотъемлемой частью. 

Корпоративные и деловые коммуникации:
\begin{enumerate}[leftmargin=0em, itemindent=3.75em, nosep, label=\arabic*)]
\item Виртуальные презентации и конференции. Инструменты (Например, PowerPoint) дополненные видео, инфографикой и анимацией, стали стандартом для донесения бизнес-идей, отчетов и стратегий. Онлайн-конференции (Например, Zoom) интегрируют видео, совместное использование экрана, интерактивные опросы и чаты, стирая географические границы и делая коммуникацию более живой;
\item Маркетинг и реклама. Современный цифровой маркетинг почти полностью построен на мультимедиа: видеоролики на YouTube и социальных сетях, интерактивные веб-сайты и лендинги, AR-приложения (дополненная реальность), позволяющие «примерить» товар, подкасты и аудиореклама. Все это направлено на создание эмоциональной связи с аудиторией и повышение вовлеченности.
\end{enumerate}

Массмедиа и социальные коммуникации:
\begin{enumerate}[leftmargin=0em, itemindent=3.75em, nosep, label=\arabic*)]
\item Конвергентные редакции. Современные новостные платформы производят контент в формате «мультимедийной истории» (digital storytelling), где один материал одновременно включает текст, фотогалерею, аудиоинтервью, видеоинфографику и интерактивные элементы, позволяя пользователю выбрать глубину погружения;
\item Социальные сети как мультимедийная среда. Платформы (Instagram, TikTok, YouTube, Facebook) изначально заточены под мультимедийный контент -- короткие видео (Reels, Shorts), прямые эфиры. Это формирует новый язык коммуникации, где визуальная и аудиальная составляющие часто первичны по~отношению к тексту.
\end{enumerate}

Социально-значимые и культурные проекты:
\begin{enumerate}[leftmargin=0em, itemindent=3.75em, nosep, label=\arabic*)]
\item Виртуальные туры и цифровые архивы. Музеи, галереи, исторические объекты используют 360° видео, интерактивные коллекции для~удаленного доступа к культурному наследию, VR-технологии, что является мощным инструментом просветительской коммуникации;
\item Публичные кампании. Социальная реклама, кампании по сбору средств или информированию общественности активно используют эмоционально заряженные видеоролики, интерактивные сайты и мобильные приложения для максимального воздействия на целевую аудиторию.
\end{enumerate}

Таким образом, мультимедийные технологии стали ключевым элементом современного образования и коммуникации. В образовательном процессе они повышают наглядность, вовлеченность и эффективность усвоения знаний, хотя их применение требует методически грамотного подхода для избежания перегрузки.
Одновременно мультимедиа формирует язык современной профессиональной и массовой коммуникации, от корпоративных презентаций до социальных сетей. Это превращает навыки создания и анализа мультимедийных сообщений в обязательную компетенцию выпускника.
Следовательно, задача образования -- не только использовать мультимедиа для обучения, но и готовить студентов к жизни в цифровом мире, где эти технологии являются основным средством для взаимодействия, познания и~творчества.
\newpage
\section{Правовые и этические аспекты работы с~мультимедийными материалами}
\subsection{Авторское право}
С увеличением цифровизации и научно-технического прогресса Интернет стал для авторов основным каналом, позволяющим им делиться своими творениями. Поскольку информация в Сети зачастую доступна свободно, включая материалы, защищенные авторским правом, любой может в любой момент легко скопировать и распространить их. Кроме того, Интернет представляет собой динамичную среду, где постоянно возникают новые процессы и явления, в то время как законодательство, с его консервативным характером, часто не успевает адаптироваться к происходящим изменениям. Таким образом, интеллектуальная собственность, особенно авторское право, оказывается в настоящее время наиболее уязвимой областью правовой защиты в цифровом пространстве.

Проблема защиты авторских прав в сети интернет долгое время не~получала внимания в России. Переломным моментом стал Федеральный закон № 149-ФЗ от 27.07.2006 «Об информации, информационных технологиях и о защите информации», который установил правила в информационной сфере и сфере технологий. Этот законодательный акт Российской Федерации определил основные понятия и термины в области правового регулирования информационных технологий, а также упорядочил отношения, возникающие при реализации прав на поиск, получение, передачу, создание и распространение информации посредством информационных технологий.

Позже приняли Федеральный закон № 187, внесший изменения в~части защиты интеллектуальных прав в информационно-телекоммуникационных сетях. Данный закон предусматривал блокировку сайтов, распространяющих нелицензионный контент, по требованию правообладателя, однако его действие распространялось только на видеоматериалы.

1 мая 2015 года вступил в силу Федеральный закон № 364, внесший поправки в закон «Об информации…» и Гражданский процессуальный кодекс РФ, расширив перечень объектов интеллектуальной собственности, включив в него книги, музыкальные произведения и программное обеспечение, но~исключив фотографии и аналогичные произведения.

Категории нарушений прав интеллектуальной собственности онлайн:
\begin{enumerate}[leftmargin=0em, itemindent=3.75em, nosep, label=\arabic*)]
\item Нарушение авторских прав, связанных с аудиовизуальными произведениями, например, незаконное распространение фильмов;
\item Нарушение авторских прав на музыкальные композиции;
\item Нарушение прав в сфере программного обеспечения, включая незаконное копирование и распространение различных программных продуктов.
\end{enumerate}

Основной метод защиты авторских прав в сети -- судебная защита. Основными судебными инстанциями, занимающимися защитой авторских прав, являются арбитражные суды, включая специализированный суд -- Суд по~интеллектуальным правам. 

Нарастающий научно-технический прогресс породил искусственный интеллект, который очень быстро овладел умением создания различных типов информации. Сейчас ИИ умеет генерировать изображения, тексты, звуки, решать различные задачи и многое другое.

Способов защиты авторских прав от ИИ существует несколько. Во-первых, технические средства защиты, такие как водяные знаки и~криптографические методы, позволяющие идентифицировать авторство человека и предотвращать несанкционированное использование его произведений ИИ. Во-вторых, можно лицензировать данные для обучения ИИ, то есть платить создателям данных за использование их материалов для обучения нейросетей, что в дальнейшем поможет избежать нарушений авторских прав при обучении ИИ.

Также запущен процесс правового регулирования данных способов. Например, в Госдуме обсуждают закон о маркировке контента, созданного нейросетями. Если его примут, то российские ИИ-сервисы обяжут маркировать созданный ими контент.
\subsection{Этика}
Этическая культура -- это доверие между создателем, пользователем и~аудиторией. Без неё медиа превращается в мусор, манипуляцию или зло.

В настоящее время существуют принципы информационной этики, но~они признаются только желающими. Четыре основных принципа присутствуют практически во всех этических кодексах: 
\begin{itemize}[label=-, leftmargin=0em, itemindent=3.35em, nosep]
\item Privacy (тайна частной жизни); 
\item Accuracy (точность в выполнении инструкций по эксплуатации систем и обработке информации);
\item Property (неприкосновенность частной собственности);
\item Accessibility (право граждан на доступность информации).
\end{itemize}

Пользователи сети Интернет должны обеспечивать уважительное и ответственное обращение с материалом, который они используют в~образовательных, исследовательских или публичных целях. Также всегда указывать автора, источник, лицензию, условия использования материала -- означает уважение к авторскому труду, что является обязательным негласным этическим правилом.

Рост цифровых технологий, социальных сетей, простота копирования -- делает нарушения авторских прав и приватности очень частыми. Появление ИИ, генеративных моделей, массового репостинга -- все это запутывает пользователей: непонятно, где автор, а где чужой контент; где согласие, а где автоматическая загрузка. Поэтому использование цифровых мультимедийных материалов требует ответственного подхода.

Работа с мультимедийными материалами -- это не просто технический процесс, но вопрос ответственности, этики и уважения. 

С появлением новых цифровых технологий (таких как искусственный интеллект) актуальны реформы законодательства об авторском праве. Изменения должны обеспечивать баланс между инновациями, защитой прав и интересами общества.
\newpage
{
\sloppy
\nocite{*}
\cleardoublepage
\addcontentsline{toc}{section}{Список использованных источников}
\bibliographystyle{plainurl}
\bibliography{report9}
}
\end{document}
